\subsection{Changing Directly from One Basis to Another}

Suppose
\[
\mathcal B=(\mathbf b_1,\mathbf b_2)
\qquad\text{and}\qquad
\mathcal C=(\mathbf c_1,\mathbf c_2)
\]
are two bases and
\[
[\mathbf u]_{\mathcal B}=(p,q).
\]

\begin{derivationbox}[Reconstruct first, then describe again]
First reconstruct the geometric vector:
\[
\mathbf u=p\mathbf b_1+q\mathbf b_2.
\]
If
\[
\mathbf b_1=(a_1,a_2),
\qquad
\mathbf b_2=(b_1,b_2),
\]
then
\[
\mathbf u=(pa_1+qb_1,\,pa_2+qb_2).
\]
Now seek $r,s$ such that
\[
\mathbf u=r\mathbf c_1+s\mathbf c_2.
\]
Writing
\[
\mathbf c_1=(d_1,d_2),
\qquad
\mathbf c_2=(e_1,e_2),
\]
we solve
\[
\begin{aligned}
 d_1r+e_1s&=pa_1+qb_1,\\
 d_2r+e_2s&=pa_2+qb_2.
\end{aligned}
\]
The solution $(r,s)$ is exactly $[\mathbf u]_{\mathcal C}$.
\end{derivationbox}

\begin{corebox}[The basis-change procedure]
A change from one basis to another consists of two familiar steps:
\[
\boxed{
\text{old coordinates}
\longrightarrow
\text{the geometric vector}
\longrightarrow
\text{new coordinates}.}
\]
At no stage does the geometric vector itself change.
\end{corebox}

\subsection{Why Perpendicular Projection Is Special}

\begin{interpretationbox}[Rectangular and oblique coordinates]
For the Cartesian basis, the coordinate directions are perpendicular unit
vectors, so coordinates can be read from perpendicular projections. For a
slanted basis this is generally false. Perpendicular shadows on
$\mathbf b_1$ and $\mathbf b_2$ need not add back to $\mathbf u$.

The correct geometric construction draws through the endpoint lines parallel to
the other basis vector. The resulting oblique parallelogram records the basis
coefficients.
\end{interpretationbox}

\begin{remarkbox}
After matrices are introduced, the same systems of equations will be written
more compactly. Matrix notation will not create a new method; it will abbreviate
the reconstruction and elimination calculations performed here.
\end{remarkbox}
